\chapter{Appendix E: Supporting SADb Records}\label{app:sadb_records}
\graphicspath{{../Figures/96-AppendixE/}{figs/96-AppendixE/}}

This appendix describes the records in the SADb Explorer export and the experimental evidence behind selected pathway labels used in Chapter~\ref{ch:database}.

\section{Record Sets and Coverage}

The SADb Explorer export contains 1,393 records across five tables. These comprise 943 Papers, 247 Models, 157 Review Papers, 42 Feedback pathways, and 4 Recorders. At the time of curation the Airtable working base had a 1,000-record limit. I retained 957 records there, comprising 561 Papers, 244 Models, 109 Review Papers, 39 Feedback pathways, and 4 Recorders. The application preserves the other 436 records as archived entries. They comprise 382 Papers and 54 records from the other tables whose paper links had been removed. The Table view includes all five record types. The Pivot and citation views use the Papers category.

Coverage of the application's Papers category is reported in Chapter~\ref{ch:database}. The following counts describe the retained maintenance subset and use a different denominator.

Of the 561 retained Papers, 523 have curation notes, 426 have preparation tags, 240 have an animal-study assessment, and 118 link to at least one named feedback pathway. Afferent classes are recorded for 86 Papers, robot or simulation translation notes for 258, and a DOI for 550. A paper can contribute to several of these counts.

Chapter~\ref{ch:database} gives the citation-graph counts. Figure~\ref{fig:sadb_clusters} uses citation proximity to locate related papers, and Figure~\ref{fig:sadb_explorer} shows the evidence recorded for a selected paper.



\section{Afferent Classes and Linked Pathways}

Table~\ref{tab:sadb_afferents} relates each receptor to a candidate model input. Table~\ref{tab:sadb_pathways} gives the pathway names as they appear in the database. In the Table view, the \emph{Ib disynaptic excitation} pathway filter returns six Papers. The record for \citet{mccrea_disynaptic_1995} carries both Ia and Ib tags and links to both \emph{Ib disynaptic excitation} and \emph{Ib disynaptic inhibition}. Chapter~\ref{ch:overall} compares the experimental conditions behind those two links.

\begin{table}[htbp]
\caption{Afferent classes, encoded signals, and pathway roles represented in the Sensory Afferent Database. References link to corpus records; Ia, II, and Ib signal summaries follow \citet{prochazka_ensemble_1998}.}
\label{tab:sadb_afferents}
\small
\begin{tabularx}{\textwidth}{@{}>{\raggedright\arraybackslash}p{0.14\textwidth}>{\raggedright\arraybackslash}p{0.16\textwidth}>{\raggedright\arraybackslash}p{0.16\textwidth}>{\raggedright\arraybackslash}X@{}}
\toprule
\textbf{Class} & \textbf{Receptor} & \textbf{Signal carried} & \textbf{Recorded pathway roles (anchors)} \\
\midrule
Ia & muscle spindle primary & muscle length and velocity & homonymous stretch excitation \citep{eccles_integrative_1958}; reciprocal inhibition of the antagonist pool \citep{jankowska_effect_1967} \\
II & muscle spindle secondary & static muscle length & midlumbar relay to motor nuclei in the cat \citep{edgley_interneuronal_1987} and evidence for a contribution to soleus activity during human stance \citep{mazzaro_afferent_2006} \\
Ib & Golgi tendon organ & whole-muscle force & extensor load feedback contributes to support during stance and to the timing of swing onset \citep{dietz_significance_2000, dominguez_candidate_2020} \\
III/IV & free nerve endings & metabolic and nociceptive signals, converging with cutaneous and group II--III afferents as flexor-reflex afferents & phase-dependent resetting of the locomotor rhythm \citep{schomburg_flexor_1998} \\
Cutaneous & skin and foot-pad receptors & contact and perturbation events & flexor excitation and phase-modulated stance correction \citep{sherrington_flexion_1910, forssberg_stumbling_1979} \\
\bottomrule
\end{tabularx}
\end{table}

\begin{sidewaystable}[p]
\caption{Selected SADb pathway records, their reported actions, and the number of Papers linked to each. Each pathway has the same count in the 943-paper export and the 561-paper retained subset.}
\label{tab:sadb_pathways}
\fontsize{10pt}{12pt}\selectfont
\setlength{\tabcolsep}{4pt}\begin{tabularx}{\linewidth}{@{}>{\raggedright\arraybackslash}p{0.18\linewidth}>{\raggedright\arraybackslash}p{0.085\linewidth}>{\raggedright\arraybackslash}X>{\centering\arraybackslash}p{0.065\linewidth}>{\raggedright\arraybackslash}p{0.23\linewidth}@{}}
\toprule
\textbf{Pathway record} & \textbf{Afferents studied} & \textbf{Reported action} & \textbf{Papers} & \textbf{Selected references} \\
\midrule
Ia monosynaptic excitation & Ia & stretch excitation of the homonymous motor pool, with the pathway gain modulated over the step cycle & 7 & \citep{eccles_integrative_1958, capaday_amplitude_1986} \\
Ia reciprocal inhibition & Ia & inhibition of the antagonist motor pool through the Ia inhibitory interneurons, rhythmically active in fictive locomotion & 6 & \citep{jankowska_effect_1967, pratt_ia_1987} \\
Ib disynaptic inhibition & I (Ia, Ib) & ankle extensor Ia and Ib afferents converge on spinal interneurons in cats. In humans, conditioning the soleus reflex from the medial gastrocnemius nerve produced less inhibition during walking than during quiet sitting & 5 & \citep{jankowska_common_1981, stephens_short_1996} \\
Ib disynaptic excitation & I (Ia, Ib) & plantaris group I input excited medial gastrocnemius motoneurons during the extensor phase of midbrain-evoked fictive locomotion in cats & 6 & \citep{mccrea_disynaptic_1995, pearson_enhancement_1998} \\
Ib stance to swing & I (includes Ib) & extensor load feedback can delay swing onset. During fictive locomotion in cats, stimulation of extensor group I afferents can prolong extension or reset the rhythm to extension & 17 & \citep{dietz_significance_2000, dominguez_candidate_2020} \\
Ia or II stance to swing & Ia, II & flexor-length gate on stance termination: lengthening of flexor muscles in late stance ends extension and resets the cycle & 10 & \citep{hiebert_contribution_1996, akay_degradation_2014} \\
type II excitatory & II & a midlumbar relay projects to motor nuclei in cats. Human ankle perturbations suggest a group II contribution to soleus activity during stance & 5 & \citep{edgley_interneuronal_1987, mazzaro_afferent_2006} \\
Cutaneous flexor excitation & Cutaneous & flexion withdrawal and the stumbling corrective reaction, phase-modulated over the cycle & 16 & \citep{sherrington_flexion_1910, forssberg_stumbling_1979} \\
\bottomrule
\end{tabularx}
\end{sidewaystable}
